%% maestro2tex -- generated Monte Carlo histogram; do not edit by hand.
\providecommand{\MaestroRoot}{include/analog/}
\begin{figure}[htbp]
  \centering
  {\pgfplotsset{compat=1.18}%
  \begin{tikzpicture}
    \begin{axis}[
      width=0.86\linewidth, height=6.0cm,
      xlabel={Phase margin~[\si{\degree}]}, ylabel={Samples},
      grid=both,
      major grid style={line width=0.2pt, draw=black!14},
      minor grid style={line width=0.2pt, draw=black!7},
      minor tick num=1,
      axis line style={line width=0.4pt, draw=black!45},
      tick style={line width=0.4pt, draw=black!45},
      tick align=outside,
      label style={font=\small}, tick label style={font=\small},
      enlarge x limits=0.02, enlarge y limits=0.10,
      every axis plot/.append style={line join=round, no marks},
      legend style={font=\footnotesize, draw=none, fill=none,
                    at={(0.5,1.02)}, anchor=south, legend columns=-1,
                    /tikz/every even column/.append style={column sep=9pt}},
      legend cell align=left,
      scaled y ticks=false, scaled x ticks=false,
      y tick label style={/pgf/number format/fixed,
                          /pgf/number format/precision=0},
      ymin=0, ymax=22.1,
      enlarge x limits=0.04, enlarge y limits=false,
      area legend
    ]
      \addplot[ybar interval, draw=mtxA, fill=mtxA!25, line width=0.7pt] table {\MaestroRoot data/tia_mc_pm560k_00.dat};
      \addlegendentry{Phase margin}
      \draw[black!70, line width=1.0pt, densely dashed] ({axis cs:45,0}|-{rel axis cs:0,0}) -- ({axis cs:45,0}|-{rel axis cs:0,0.90})
        node[above right, font=\footnotesize, inner sep=1pt] {floor \SI{45}{\degree}};
    \end{axis}
  \end{tikzpicture}}
  \caption[{Transimpedance loop phase margin at maximum gain over 179 converged samples of 200 (process and mismatch),\,\dots}]{Transimpedance loop phase margin at maximum gain over 179 converged samples of 200 (process and mismatch), \SI{37}{\celsius}. Every converged sample clears the \SI{45}{\degree} floor, the worst at \SI{46.6}{\degree}, which is what validates $C_f = \SI{22}{\pico\farad}$ --- the corner run behind Table~\ref{tab:tia-cl-stability} gave \SIrange{61.2}{64.9}{\degree} and carries no mismatch at all. The distribution is right-skewed, so \(p_1\)/\(p_{99}\) and worst case are quoted rather than mean \(\pm\,3\sigma\). 48 of 400 samples did not converge and are excluded (Section~\ref{sss:tia-mc}).}
  \label{fig:tia-mc-pm560k}
\end{figure}
